% !TEX root = ../diplomarbeit.tex

\chapter{Einleitung\authorA}

In diesem Kapitel werden die Ziele dieser Diplomarbeit festgehalten. Die Gliederung der Arbeit wird durch einen kurzen Überblick über die Kapitel beschrieben.

\section{Motivation, Problemstellung und Zielsetzung}

% Dieser Abschnitt erläutert die Motivation für die Themenwahl, beschreibt das zu lösende Problem und definiert die konkreten Ziele der Arbeit. Eine klare Problemstellung und explizite Zielsetzung sind wesentlich, um den Rahmen der Arbeit abzustecken und die Erfolgskriterien zu definieren.

\placeholderbox{Beispiel}{%
\subsection*{Motivation}
In der modernen Softwareentwicklung stellt die Verwaltung von Abhängigkeiten in großen Projekten eine zunehmende Herausforderung dar. Entwicklerteams verbringen durchschnittlich 20\,\% ihrer Arbeitszeit mit der Analyse und Behebung von Versionskonflikten zwischen Bibliotheken.

\subsection*{Problemstellung}
Bestehende Dependency-Management-Tools bieten zwar automatische Aktualisierungen an, berücksichtigen jedoch nicht die semantische Kompatibilität auf API-Ebene. Dies führt zu Breaking Changes, die erst zur Laufzeit erkennbar werden und kostspielige Hotfixes erfordern.

\subsection*{Zielsetzung}
Ziel dieser Arbeit ist die Entwicklung eines intelligenten Dependency-Analyzers, der mittels statischer Code-Analyse die tatsächliche API-Nutzung erfasst und automatisch kompatible Versionen vorschlägt. Das System soll mindestens 90\,\% der Breaking Changes bereits vor der Integration erkennen und die Analysezeit unter 5 Minuten für typische Projekte halten.
}




\section{Gliederung der Arbeit}

% Dieser Abschnitt gibt einen Überblick über den Aufbau der Arbeit und beschreibt kurz den Inhalt der einzelnen Kapitel. Dies ermöglicht Lesenden eine schnelle Orientierung und zeigt die logische Struktur der Argumentation auf.

\placeholderbox{Beispiel}{%
Die vorliegende Arbeit gliedert sich in sieben Kapitel:\\

\textbf{Kapitel 1 -- Einleitung} führt in die Thematik ein, beschreibt die Problemstellung und definiert die Zielsetzung der Arbeit.

\textbf{Kapitel 2 -- Grundlagen} vermittelt das notwendige theoretische Fundament, definiert zentrale Begriffe und stellt den aktuellen Stand der Technik dar.

\textbf{Kapitel 3 -- Konzepte} präsentiert die entwickelten Lösungsansätze, beschreibt die Systemarchitektur und begründet getroffene Designentscheidungen.

\textbf{Kapitel 4 -- Implementierung} dokumentiert die technische Umsetzung, erläutert kritische Implementierungsdetails und beschreibt verwendete Algorithmen.

\textbf{Kapitel 5 -- Prototyp und Evaluation} stellt den entwickelten Prototyp vor, beschreibt durchgeführte Tests und evaluiert die Ergebnisse hinsichtlich der definierten Zielsetzung.

\textbf{Kapitel 6 -- Zusammenfassung und Ausblick} fasst die wichtigsten Erkenntnisse zusammen, reflektiert kritisch die Ergebnisse und skizziert mögliche Weiterentwicklungen.

\textbf{Kapitel 7 -- Projektmanagement} dokumentiert die Projektorganisation, Aufgabenverteilung und den zeitlichen Ablauf der Diplomarbeit.
}